\documentclass{article}
\usepackage{ctex}
\usepackage{amsmath}
\usepackage[paperwidth=19.6cm,paperheight=7.55cm,margin=0.15cm]{geometry}
\usepackage{tikz}
\usetikzlibrary{arrows.meta,calc,decorations.pathreplacing,positioning}
\setCJKmainfont{Songti SC}
\definecolor{temp}{HTML}{B64342}
\definecolor{moist}{HTML}{0F4D92}
\definecolor{radiusc}{HTML}{42949E}
\definecolor{neutral}{HTML}{767676}
\definecolor{pale}{HTML}{F4F5F7}
\definecolor{palered}{HTML}{F8E9E7}
\definecolor{paleblue}{HTML}{E9F0F8}
\tikzset{
  every node/.style={font=\small},
  flow/.style={-{Latex[length=2.2mm,width=1.5mm]},line width=.8pt},
  dim/.style={{Latex[length=1.8mm,width=1.2mm]}-{Latex[length=1.8mm,width=1.2mm]},line width=.65pt},
  note/.style={rounded corners=2pt,draw=neutral!70,fill=pale,inner sep=4pt,align=left},
  panel/.style={font=\bfseries\large,anchor=west}
}
\begin{document}
\thispagestyle{empty}
\noindent
\begin{tikzpicture}[x=1cm,y=1cm]
  \node[panel] at (0,6.5) {(a) 圆柱几何与端面口径};
  \fill[pale] (0.7,1.55) rectangle (8.0,5.15);
  \draw[line width=.9pt] (0.7,1.55) rectangle (8.0,5.15);
  \draw[dashed,neutral] (4.35,1.4) -- (4.35,5.3);
  \draw[dash pattern=on 2pt off 1.5pt,neutral] (0.7,3.45) -- (8.0,3.45);
  \node[above,neutral] at (4.35,5.2) {几何中截面 $z=0$};
  \node[above,neutral] at (1.45,3.45) {轴线 $r=0$};
  \draw[dim] (0.7,1.05) -- node[fill=white,inner sep=1pt] {$L=25\,\mathrm{cm}$} (8.0,1.05);
  \draw[dim] (8.45,3.45) -- node[right,fill=white,inner sep=1pt] {$R(t)$} (8.45,5.15);
  \draw[flow,temp] (1.8,5.72) -- (1.8,5.18);
  \draw[flow,temp] (2.7,5.72) -- (2.7,5.18);
  \draw[flow,moist] (5.6,5.72) -- (5.6,5.18);
  \draw[flow,moist] (6.5,5.72) -- (6.5,5.18);
  \node[temp,above] at (2.25,5.72) {$h(T_{\rm air}-T_s)$};
  \node[moist,above] at (6.05,5.72) {$h_m(C_s-C_{\rm eq})$};
  \draw[flow,temp] (0.05,4.35) -- (0.68,4.35);
  \draw[flow,moist] (0.05,2.55) -- (0.68,2.55);
  \node[left,temp] at (0.02,4.35) {端面换热};
  \node[left,moist] at (0.02,2.55) {端面传质};
  \node[note,text width=3.2cm] at (6.0,2.25) {一维主模型：端面绝热且不透湿；\\二维审查：端面与侧面同样暴露。};

  \node[panel] at (9.2,6.5) {(b) 径向模型与收缩坐标};
  \coordinate (O1) at (11.2,4.15);
  \coordinate (O2) at (15.7,4.15);
  \fill[paleblue] (O1) circle (1.55);
  \draw[moist,line width=1pt] (O1) circle (1.55);
  \fill[palered] (O2) circle (1.00);
  \draw[temp,line width=1pt] (O2) circle (1.00);
  \draw[flow,radiusc] (O1) -- ++(1.55,0) node[midway,above] {$r$};
  \draw[flow,radiusc] (O2) -- ++(1.00,0) node[midway,above] {$r$};
  \draw[flow,neutral] (12.95,4.15) -- node[above] {$R_0\rightarrow R(t)$} (14.5,4.15);
  \node[below=1.75cm of O1] {固定域：$0\le r\le R_0$};
  \node[below=1.2cm of O2,align=center] {移动域：$\xi=r/R(t)$\\$0\le\xi\le1$};
  \node[note,text width=7.0cm] at (13.45,1.15) {$Q1$--$Q3$：$R(t)=R_0=2\,\mathrm{cm}$；\\$Q4$：仅径向仿射收缩、长度固定，且材料速度与网格速度一致。};
  \draw[decorate,decoration={brace,amplitude=5pt,mirror},line width=.65pt,neutral]
    (9.55,0.35) -- (17.45,0.35) node[midway,below=5pt,align=center] {轴线取对称零通量；侧面采用对流换热与传质边界};
\end{tikzpicture}
\end{document}
